\chapter{Описание стенда и подготовка}
\label{ch:stand}

\section{Конфигурация виртуального стенда}

Лабораторная работа выполняется на виртуальной машине \texttt{gateway} под управлением ОС \textbf{Ubuntu 20.04.5 LTS (Focal Fossa)}, развёрнутой в среде Oracle VirtualBox. ВМ имеет два сетевых адаптера:
\begin{itemize}
    \item \texttt{enp0s3} --- WAN-интерфейс, подключён в режиме NAT; получает адрес \texttt{10.0.2.15/24}, шлюз \texttt{10.0.2.2};
    \item \texttt{enp0s8} --- LAN-интерфейс, подключён к внутренней сети (Internal Network); статический адрес \texttt{192.168.N.1/24}.
\end{itemize}

Лаба 5 выполняется на той же ВМ, что и лаба 4. К моменту начала лабы 5 уже настроены: маршрутизация (ip\_forward + iptables MASQUERADE) и DHCP-сервер \texttt{isc-dhcp-server}.

\section{Предварительная подготовка}

Перед установкой BIND9 необходимо выполнить ряд подготовительных шагов.

\subsection{Установка hostname}

Имя хоста сервера задаётся командой:
\begin{lstlisting}
hostnamectl set-hostname SERVERHOSTNAME
\end{lstlisting}
где \texttt{SERVERHOSTNAME} --- выбранное имя ВМ (например, \texttt{gateway}). Файл \texttt{/etc/hosts} редактируется вручную:
\begin{lstlisting}
nano /etc/hosts
# добавить строку:
127.0.1.1   SERVERHOSTNAME
\end{lstlisting}

\subsection{Удаление DNS DNAT-правил}

В лабе 4 в файл \texttt{/etc/iptables/rules.v4} были добавлены правила перенаправления DNS-запросов на Google~8.8.8.8. Теперь эти правила необходимо удалить, поскольку DNS-сервер будет обрабатывать запросы самостоятельно:
\begin{lstlisting}
nano /etc/iptables/rules.v4
# удалить строки:
-A PREROUTING -i enp0s8 -p tcp -m tcp --dport 53 \
  -j DNAT --to-destination 8.8.8.8:53
-A PREROUTING -i enp0s8 -p udp -m udp --dport 53 \
  -j DNAT --to-destination 8.8.8.8:53
\end{lstlisting}

После редактирования применить изменения и перезагрузить ВМ:
\begin{lstlisting}
reboot
\end{lstlisting}

\section{Установка BIND9}

Установка пакетов производится стандартными средствами менеджера пакетов:
\begin{lstlisting}
apt-get update
apt install bind9
apt install dnsutils
\end{lstlisting}

Пакет \texttt{dnsutils} предоставляет утилиты \texttt{nslookup} и \texttt{dig}, необходимые для проверки работы DNS.

\endinput
